% Finite coefficient realization proof. An executed Q8 table follows in
% sections/D_final_table.tex.

\subsection{A finite realization calculus for the alternating projector}

\begin{lemma}[Finite coefficient data may be tensored with a finite category]
\label{D:finite-coefficient-tensor}
Let $K$ be a finite connected polyhedron.  Consider a finite direct
homotopy-coherent diagram of finite polyhedra under and over $K$,
including only finitely many specified homotopies and higher homotopies.
Here ``over $K$'' is derived: a map $f:Y\to Z$ is accompanied by a
specified homotopy $r_Zf\simeq r_Y$ relative to the section $K$.
Let $A$ be the classifying space of a finite pointed category.
The diagram obtained by tensoring these retractive spaces with $A$,
and then applying finitely many products and finite homotopy colimits,
has a model built from finite categories and functors.  The specified
homotopies have finite categorical representatives in this model.
The assertion also allows a finite list of functors, natural
transformations, and finite Grothendieck constructions in the variable
pointed category $A$.
\end{lemma}
\begin{proof}
We describe all finite constructions needed.  A single map with its
specified over-$K$ homotopy already admits a finite strict model.
If $f:Y\to Z$ fixes $K$ and $H:r_Zf\simeq r_Y$ is the given
homotopy relative to $K$, form its relative mapping cylinder
\[
 Z_f=(Y\times[0,1])\cup_f Z\,/\,((k,t)\sim k\text{ for }k\in K).
\]
Choose the orientation of $H$ so that it is $r_Y$ at the source
end and $r_Zf$ at the target end.  It defines a projection
$Z_f\to K$, agreeing with $r_Z$ on $Z$.  The two inclusions
$Y\to Z_f\leftarrow Z$ are literal maps under and over $K$,
and the second is a weak equivalence on underlying spaces.  It is
therefore invertible in the derived retractive category; the ordinary
cylinder contraction need not preserve the projection to $K$.
Thus they represent the given derived map.  A homotopy between
such data is treated by the same relative mapping-cylinder
construction with an additional interval parameter; a specified
$j$-fold homotopy uses a $j$-cube.  Every space used is finite.
No path-space replacement is being asserted finite.

Here is a relative finite-model construction which avoids any assertion
that an arbitrary finite family of maps admits one simultaneous
triangulation.  First make the finitely many coefficient maps and
parameter homotopies piecewise linear.  Proceed in increasing parameter
dimension: the boundary data have already been made piecewise linear,
the homotopy extension property corrects the old filling to that
boundary, and relative simplicial approximation supplies a piecewise
linear filling.  Only finitely many fillings are used.

Retain a map as the two inclusions into its relative mapping cylinder,
as in the first paragraph, instead of triangulating the map itself.
Treat a specified homotopy by its relative prism and mapping cylinder,
and do the same for the finitely many higher parameter cubes.  Duplicate
each occurrence of an object before this construction.  Attach these
cylinders and prisms in the order of the finite direct pasting shape.
Each old occurrence is a subpolyhedron of the resulting finite
polyhedron: attaching a cylinder or a parameter prism does not identify
distinct points of the old target.  All copies of the specified section
are identified with the one copy of $K$.  The given projection
homotopies agree on the attaching boundaries, so they define one
piecewise linear retraction $r:Z\to K$ on the finite union $Z$.
There are finitely many distinguished subpolyhedra of $Z$, namely the
object occurrences, cylinders, prisms, and their faces.  Their maps are
inclusions.  The target-to-cylinder inclusions which are inverted are
underlying homotopy equivalences.  Thus this finite diagram of
inclusion roofs represents precisely the original finite pasting data,
including the specified homotopies.  This is the cylinder description
of finite bar rectification; finiteness follows also by indexing its
cubes by strings of nonidentity arrows in the finite direct shape.

We record why the single retraction causes no hidden subdivision
cycle.  Start with finite polyhedral decompositions of $Z$ and $K$
on which $r$ is affine, preserving the section and every distinguished
subpolyhedron.  Subdivide $K$ by the images under $r$ of all faces of
the decomposition of $Z$, using the finite arrangement of their
supporting hyperplanes.  Subdivide each old cell of $Z$ by inverse
images of these target cells.  Every resulting cell maps onto one
target cell: its faces are intersections of an old face with the
inverse image of a target face, and the target decomposition was
chosen to subdivide the images of all old faces.  This gives finite
polyhedral face posets $P_Z,P_K$ and an order-preserving map
$P_r:P_Z\to P_K$ taking a cell to its image cell.  On the section
it is the identity.  Every distinguished inclusion commutes strictly
with $P_r$.

These poset models can even be compared over $K$ literally.  Choose
an interior point of each target cell.  In each source cell choose an
interior point over the chosen point of its image cell; surjectivity
onto that cell guarantees such a point.  Use the same choices on all
distinguished subpolyhedra, and on the section use the target choices.
Coning the boundary triangulations to these interior points gives the
usual face-poset homeomorphisms.  Since $r$ is affine on each source
cell, they commute with $r$, and all inclusion comparisons commute
as well.  Consequently the face-poset inclusion roofs are a finite
strict diagram of retractive categories representing the given data.

No self-idempotent on a fixed finite polyhedron has been strictified.
Repeated occurrences are connected by the displayed finite equivalence
roofs.  In particular the retraction triangle is treated as relative
structure on the one map $Z\to K$, not incorrectly declared to be a
direct indexing triangle.

Here is the tensor construction itself.  For a strict retractive
coefficient object $K\xrightarrow{i}Y\xrightarrow{r}K$, $ri=1$,
and a pointed category $C$ with $BC=A$, first take the finite
Grothendieck construction $D$ of
\[
 K\times C\ \longleftarrow\ K\ \longrightarrow\ Y.
\]
The first map uses the base object of $C$.  There are strict functors
$D\to Y\times C$ and $D\to K$, induced respectively by $i$ and
by $r$.  The Grothendieck construction of
\[
 Y\times C\ \longleftarrow\ D\ \longrightarrow\ K
\]
models the derived pushout
\[
 (Y\times A)\ \cup^h_{(K\times A)\cup^h_K Y}\ K.
\]
This is the tensor of $Y$ with $A$ in retractive spaces over $K$.
Every category and every hom-set is finite.  The recipe is strictly
functorial for strict coefficient maps and for functors of pointed
categories.  Products with the finite interval and cubical
coefficient categories carry the specified homotopies along as
well.  The homotopy-colimit theorem transfers the preceding
coefficient equivalences through these constructions.  Additional
finite Grothendieck constructions remain finite, and their
comparison maps and natural transformations are the same
explicit tuple formulas.  This proves the lemma.
\end{proof}

\begin{lemma}[The coefficient matrices and their controlled homotopies]
\label{D:coefficient-matrices}
Let $K$ be a finite connected CW complex of dimension $k$, with
fundamental group $\pi$, and take $m>k+1$, $m\geq3$.
Put $Y=K\vee\bigvee_{j=1}^N S^m$, as a retractive space over $K$.
Every endomorphism of the free $\mathbb Z\pi$-module on the sphere
cells is represented by a map $Y\to Y$ in the derived retractive
homotopy category.  In particular, for $q=1-e$, choose coefficient maps representing
$e$ and $q$ and one derived nullhomotopy for each of $eq=0$ and
$qe=0$.  These two nullhomotopies have finite coefficient
representatives, including their specified over-$K$ homotopies.
Their cellular chain homotopies in degree $m$ can be taken zero.
\end{lemma}
\begin{proof}
The retraction induces a split surjection on every homotopy group.
Its homotopy fibre is simply connected, and the relative Hurewicz
isomorphism for the $m$-cell attachments identifies its $m$-th
homotopy group with the free $\mathbb Z\pi$-module on the attached
sphere cells.  A map under and derived-over $K$ is specified on
these finitely many sphere cells by elements of that group.
Composition acts on their coefficients by matrix multiplication,
with the chosen left/right module convention.  Thus choose maps representing $e,q$ and two paths in the derived
mapping space representing $eq\simeq0$ and $qe\simeq0$.  Each such
path is an ordinary homotopy of the underlying maps accompanied
by its two-parameter over-$K$ coherence.  The source of each
adjoint map is a finite CW complex times finitely many intervals,
so these are finite coefficient data.  We do not assert that an
arbitrary list of equalities in a homotopy category admits
arbitrarily prescribed compatible higher homotopies.

More concretely, write each matrix entry as a finite integral sum
of elements of $\pi$.  Pinch an $m$-sphere into the required number
of summands, use a reflection for a negative coefficient, and
whisker each summand along a finite edge path in $K$ representing
the group element.  The over-$K$ homotopy contracts the path and
its reverse after collapsing the sphere.  The relative co-H
operations are abelian because $m\geq3$; reassociation,
commutation, inverse cancellation and finitely many path
homotopies in $K$ give representatives of the two chosen
nullhomotopies.  No relation between these two nullhomotopies
is required: they occur on different parity summands.
These are finite parameter homotopies, to which
Lemma~\ref{D:finite-coefficient-tensor} applies.

The coefficient complex $Y$ has no cells above dimension $m$.
A cellular homotopy from a sphere map to another sphere map
therefore has zero cellular chain homotopy out of degree $m$.
There are no intervening base cells in degrees $m-1$ or $m$,
by the choice $m>k+1$.  This gives the claimed chain control.
Tensoring such a homotopy with a parity factor preserves the
zero relative chain homotopy, since it acts on the coefficient
factor and is the identity on the parity factor.
\end{proof}

\begin{theorem}[Finite-category realization of projective finiteness obstructions]
\label{D:Wall-realization}
Let $\pi$ be a finitely presented group and let
$e:(\mathbb Z\pi)^N\to(\mathbb Z\pi)^N$ be an idempotent
endomorphism.  There is a finite connected category $\mathcal C_e$
whose fundamental group is $\pi$, whose classifying space is
finitely dominated, and whose Wall obstruction is
\[
 w(B\mathcal C_e)=(-1)^{m+2}[(1-e)(\mathbb Z\pi)^N]
 \quad\text{in }\widetilde K_0(\mathbb Z\pi),
\]
where $m$ may be any sufficiently large chosen integer.
In particular D1 has a negative answer, including for fundamental
groups $Q_8$ and $C_{23}$, by the explicit projectives of
Section~\ref{D:explicit-projectives}.
\end{theorem}
\begin{proof}
Choose a finite connected presentation complex $K$ with group $\pi$
and choose $m$ as in Lemma~\ref{D:coefficient-matrices}.
Use the split parity space
$A=BM_0\vee BM_1$, with the finite categorical wedge model of
Theorem~\ref{D:parity-monoids}.  Tensor the coefficient object
$Y=K\vee\bigvee^N S^m$ with $A$ in retractive spaces over $K$.
Its underlying space has type
\[
 V=K\vee\bigvee_{j=1}^N(S^m\wedge A).
\]
Its relative universal-cover chains are homotopy equivalent to
one free $(\mathbb Z\pi)^N$ in every degree at least
$d=m+2$, with zero differential.  Choose the basis supplied by
the primitive spherical generators of the parity monoids.

Let $f:V\to V$ have coefficient $e$ on the $M_0$ summands and
$q=1-e$ on the $M_1$ summands.  Let $f^c$ use $q$ on the first
parity and $e$ on the second.  These are built from the same two
chosen coefficient maps, with their positions exchanged.  It is a map in the derived retractive
category, represented by finite coefficient data tensored with
the parity categories.  The uniform shift maps of
Theorem~\ref{D:parity-shifts} give
\[
 \sigma:\Sigma_K V\longrightarrow V.
\]
They exchange the parity summands, sending each suspended chosen
generator to a unit $\pm1$ times the next generator. No common
suspension-sign convention is needed below. The shifts act on the
parity factor; the coefficient maps act on $Y$.

Form the relative mapping cone
$C=\operatorname{Cof}_K(f)$, and let $j:V\to C$ denote its
inclusion.  The finite coefficient nullhomotopy
$ff^c\simeq0$ defines a map
\[
 \ell:\Sigma_K V\longrightarrow C
\]
whose composite with the cone quotient is $\Sigma_Kf^c$.
One can construct it directly on the two suspension cones:
on one cone use $f^c$ into the source cone of $C$, and on the
other use the specified nullhomotopy.  Thus no stable exactness
assertion is needed to produce this unstable map.
Put
\[
 g=\ell+jf^c\sigma:\Sigma_KV\longrightarrow C,
 \qquad T=\operatorname{Cof}_K(g).
\]
The sum is the relative co-H sum. The explicit table below instead
uses a pinch with branch degrees $(1,-1)$; this only changes a unit
sign in the following calculation. Every construction here is
a finite product, finite homotopy colimit, finite coefficient
homotopy, or one of the explicitly represented parity shifts.
Lemma~\ref{D:finite-coefficient-tensor} therefore supplies a
finite category $\mathcal C_e$ and a homotopy equivalence
$B\mathcal C_e\simeq T$.

We calculate the chains, retaining the specified homotopy.
Write $E=(\mathbb Z\pi)^N$ and renumber the bottom sphere degree
$d$ as zero.  On the relative formal chain model the components
of $f$ alternate $e,1-e$.  The cone $C/K$ has the mapping-cone
complex $K$ of Proposition~\ref{D:alternating-projector}.
The lower coordinate of $\ell$ in degree $n>0$ is
$1-f_{n-1}=f_n$.  Its upper chain-homotopy coordinate is zero
by the finite coefficient control in
Lemma~\ref{D:coefficient-matrices}. The summand $jf^c\sigma$ has
upper coordinate $\epsilon_n(1-f_n)$ and lower coordinate zero,
where $\epsilon_n\in\{1,-1\}$. Hence $g$ has the form
\[
 x\longmapsto(\epsilon_n(1-f_n)x,f_nx).
\]
This is an isomorphism from $E$ to
$(1-f_n)E\oplus f_nE=H_n(C/K)$, with inverse
$(u,v)\mapsto\epsilon_nu+v$. Thus the proof of
Proposition~\ref{D:alternating-projector}, with this signed
coordinate, shows that $C_*(\widetilde T,\widetilde K)$ is
chain-homotopy equivalent to $(1-e)E[d]$.

All spaces in the construction have the same fundamental group
as $K$: their newly attached factors are simply connected, and
the two relative cones apply van Kampen to maps inducing the
identity on $\pi$.  Thus $\pi_1(T)=\pi$.
The universal-cover complex of $K$ is finite free, while the
relative complex just computed is perfect.  The exact triangle
of the pair therefore makes $C_*(\widetilde T)$ chain-homotopy
equivalent to a bounded complex of finitely generated projectives.
Together with the finite presentation of $\pi$, Wall's finite-
domination criterion \cite{wall1965D} makes $T$ finitely dominated. The
additivity of the projective Euler class gives exactly
$(-1)^d[(1-e)E]$, since the finite base contributes zero in
reduced $K_0$.  A finitely dominated CW complex has a finite-
dimensional CW model, for example the mapping telescope of a
homotopy idempotent on a finite dominating complex.  This is the
required finiteness property.
\end{proof}
